\begin{table}[H]
    \caption{Results on Canadian insurance dataset. Reported performance metrics are averaged over $30$ runs. }
    %For both statistics, the higher is better. On the other hand, overfitting can be assessed by the difference in magnitudes between training set and out-of-sample set.
    \centering
    \begin{tabular}{|p{3.3cm}||p{2.1cm}|p{1.8cm}|p{1.8cm}|}
         \hline
         Model &  $\qadj~(\mathrm{S.E.})$ & $\hat{v}~(\mathrm{S.E.})$ & $\hat{\mathcal{R}}~(\mathrm{S.E.})$\\
         \hline
    %TRAIN
         \multicolumn{4}{|c|}{Training Set ($n=21,000$, $p=40$)} \\
         \hline
         Qini-based UR  &  & & \\
         Causal Forest  &  & & \\
         Causal Forest (H) &  & & \\
         Random Forest (ED) &  & & \\
         Random Forest (KL) &  & & \\
         SMITE (TO) &  & & \\
         SMITE (IE) &  & & \\
         \hline
    %VALID
         \multicolumn{4}{|c|}{Validation Set ($n=14,000$, $p=40$)} \\
         \hline
         Qini-based UR  &  & & \\
         Causal Forest  &  & & \\
         Causal Forest (H) &  & & \\
         Random Forest (ED) &  & & \\
         Random Forest (KL) &  & & \\
         SMITE (TO) &  & & \\
         SMITE (IE) &  & & \\
         \hline
    %TEST
         \multicolumn{4}{|c|}{Out-of-sample ($n=15,000$, $p=40$)} \\
         \hline
         Qini-based UR  &  & & \\
         Causal Forest  &  & & \\
         Causal Forest (H) &  & & \\
         Random Forest (ED) &  & & \\
         Random Forest (KL) &  & & \\
         SMITE (TO) &  & & \\
         SMITE (IE) &  & & \\
         \hline
    \end{tabular}
        \label{tab:realdata_perf}
\end{table}